\section{Ý nghĩa cho thiết kế thực nghiệm}
\label{sec:implications}

Kết quả đo đặc trưng giúp định hướng thiết kế baseline ở WEEK03 theo giả thuyết RQ1:
chuỗi có \(F_S\)/\(F_T\) cao thì kỳ vọng DLinear tốt hơn; chuỗi không dừng hoặc
hay dịch mức thì NLinear có lợi thế hoặc hai mô hình cho kết quả tương đương.

\begin{examplebox}{Phân nhóm đối chứng}
  \textbf{Nhóm nghiêng DLinear (mùa vụ/xu hướng rõ):}
  Weather, ETTh2, Electricity --- \(F_S\) cao, chuỗi dừng.\\[0.3em]
  \textbf{Nhóm nghiêng NLinear (không dừng, không mùa vụ):}
  Exchange-Rate, VIC --- \(F_S = 0\), không bác bỏ unit root.\\[0.3em]
  \textbf{Trường hợp trung gian:} ETTh1 --- xu hướng mạnh (\(F_T\) cao)
  nhưng mùa vụ yếu; sẽ báo cáo riêng, không dùng đơn lẻ để kết luận H1.
\end{examplebox}

\begin{enumerate}
  \item Khi vẽ biểu đồ scatter cho RQ1, trục \(X\) ưu tiên dùng \(F_S\)
        (vì phân tách rõ nhất trong bảng); \(F_T\) và \(p_{\mathrm{ADF}}\)
        có thể đưa vào trục phụ hoặc ghi chú.
  \item VIC (\(n \approx 3{,}800\)) và Exchange (\(n \approx 7{,}600\))
        ngắn hơn nhiều so với Weather (\(n \approx 53{,}000\)).
        Cần ghi rõ lưới lookback/horizon và số seed trong protocol
        để tránh kết luận quá mức từ mẫu nhỏ.
  \item Bảng~\ref{tab:profiles} là đầu vào cố định cho heatmap WEEK04.
        Nếu thay đổi chu kỳ \(m\) của STL, cần chạy lại script và cập nhật CSV.
\end{enumerate}

\vspace{0.5em}
\noindent\textit{File liên quan:}
\code{Working_Files/TA_reports/WEEK02_Data_Profiles.csv},
\code{Working_Files/src/data_chars.py}.
